\section{First-pass costings}\label{sec:results}

\subsection{Eligible flows}

Over the four quarters to June 2026, ONS redundancy data give an annual
flow of \TaaFlowManuf\ manufacturing redundancies and \TaaFlowAll\ across
all industries; the trade-exposed memo population is \TaaFlowExposed. One
caveat travels with the narrow flow: the final quarter's manufacturing
figure (\TaaRedManufAprJunTwentySix) is anomalously low against the three
preceding quarters (\TaaRedManufJulSepTwentyFive,
\TaaRedManufOctDecTwentyFive, \TaaRedManufJanMarTwentySix)---LFS
redundancy cells are small and noisy---so the narrow costings likely sit
toward the low end of their steady-state range.

\subsection{Annual costs}

Table~\ref{tab:taa-costs} reports each arm's gross annual cost at central
take-up under both eligibility rules, with the trade-exposed memo
column. Every cell is generated by the pipeline from the pinned inputs;
low--high scenario totals bracket the central column.

\begin{table}[htbp]
\centering
\caption{Gross annual cost by arm and eligibility, central take-up
(\pounds m/year)}
\label{tab:taa-costs}
\small
\begin{tabular}{lrrr}
\toprule
Arm & Narrow & Broad & Memo: trade-exposed \\
 & (manufacturing) & (all industries) & (30\% of manufacturing) \\
\midrule
A. Wage insurance & \TaaWageInsCostCentralNarrow & \TaaWageInsCostCentralBroad & \TaaWageInsCostCentralExposed \\
B. UI extension & \TaaUiExtCostCentralNarrow & \TaaUiExtCostCentralBroad & \TaaUiExtCostCentralExposed \\
C. Retraining stipend & \TaaTrainCostCentralNarrow & \TaaTrainCostCentralBroad & \TaaTrainCostCentralExposed \\
D. UC disregard (exposure bound) & \TaaUcDisCostCentralNarrow & \TaaUcDisCostCentralBroad & \TaaUcDisCostCentralExposed \\
\midrule
Total & \TaaTotalCostCentralNarrow & \TaaTotalCostCentralBroad & \TaaTotalCostCentralExposed \\
Scenario bracket (low--high) & \TaaTotalCostLowNarrow--\TaaTotalCostHighNarrow & \TaaTotalCostLowBroad--\TaaTotalCostHighBroad & \TaaTotalCostLowExposed--\TaaTotalCostHighExposed \\
\bottomrule
\end{tabular}
\begin{minipage}{.94\textwidth}\footnotesize
\emph{Notes:} gross steady-state envelopes; no tax or Universal Credit
clawback, no behavioural response, arms costed independently and summed.
Arm D is a labelled maximum-exposure bound, not a costing
(Section~\ref{sec:design}). Scenario brackets move all take-up,
re-employment and penalty dials together in the cost-increasing
direction, so they are outer bounds, not confidence intervals.
\end{minipage}
\end{table}

Per-beneficiary magnitudes make the arms comparable to existing
benefits. The central wage-insurance payment is
\pounds\TaaWageInsPerWorkerCentralManuf\ a year for a re-employed
manufacturing worker (half of a 21.5 per cent earnings loss on the
\pounds\TaaMeanWeeklyManuf\ mean full-time manufacturing wage)---roughly
the annual value of the proposed Unemployment Insurance itself. The UI
extension is worth \pounds\TaaUiExtPerWorkerCentral\ per claimant (the
gap between \pounds\TaaUiWeekly\ a week for nine months and the current
\pounds\TaaJsaWeekly\ New Style JSA for six). The retraining stipend is
\pounds\TaaTrainPerWorker\ flat; the UC bound prices six months of the
single-adult standard allowance for capital-gated claimants.

\subsection{Scale: the ad-hoc comparison}

Government money committed around the two current steel interventions is
roughly \pounds\TaaSteelSupportTotal\,million
(\pounds\TaaSteelPortTalbot\,million at Port Talbot,
\pounds\TaaSteelScunthorpe\,million at Scunthorpe to January 2026, the
latter still accruing). The central narrow-eligibility programme---all
four arms, every manufacturing redundancy in the country---costs
\TaaTotalCentralPctSteelNarrow\ per cent of that; the trade-exposed memo
programme costs \TaaTotalCentralPctSteelExposed\ per cent. A designed,
national, rules-based insurance system for manufacturing displacement is,
in gross annual cost, the same kind of money the UK has committed to two
plants in two years. The broad programme is different money
(\TaaTotalCentralPctSteelBroad\ per cent of the steel anchor): extending
trade-adjustment instruments into general redundancy insurance is a
welfare-state reform, not a trade policy, and the bracket between the
narrow and broad columns is the price of avoiding a certification
mechanism.

The comparison cuts both ways and should be read symmetrically: the
steel money buys plant continuity, supply-chain security and local
transition support, not only worker insurance; and the insurance
programme would cover workers the ad-hoc packages never reach. What the
juxtaposition establishes is narrower and robust to the caveats---the
fiscal scale of a statutory instrument is not what has prevented one.
